
\begin{abstract}[it]
El trabajo propone una metodología para estimar la tasa de desempleo mediante técnicas de Estimación de Áreas Pequeñas (SAE) basadas en modelos mixtos con componentes espaciales, en los municipios de los departamentos del  Cauca y Valle del Cauca, Colombia.  Las fuentes oficiales, como la Gran Encuesta Integrada de Hogares (GEIH) y la Encuesta Multipropósito (EM), tienen diseños muestrales que no garantizan estimaciones confiables en áreas pequeñas (como municipios o Unidades de Planeamiento Zonal – (UPZ), lo que resulta en estimadores directos con alta varianza y baja precisión estadística. Esta limitación técnica se complementa con la escasez de información censal pública actualizada y la baja disponibilidad de registros administrativos que podrían servir como datos auxiliares. La metodología propuesta puede servir como insumo para obtener estimaciones del desempleo municipal más preciso y replicable que permita a los gobiernos territoriales y entidades nacionales fortalecer la formulación de políticas públicas focalizadas, diseñar estrategias más efectivas de desarrollo económico y reducción de la pobreza, y orientar mejor los recursos públicos hacia las zonas con mayores necesidades. Además, metodología  puede servir de referencia para la extensión a la totalidad de los municipios de Colombia y la automatización del proceso de estimación para un monitoreo continuo del mercado laboral.
\end{abstract}

%-XXXXXXXXXXXXXXXXXXXXXXXXXXXXXXXXXXXXXXXXXXXXXXXXXXXXXXXXXXXX
% Este resumen hay que actualizarlo con detalels de la metodología
% y resultados.
% Adicionalmente, hay que escribirlo en presente y no en futuro,
% pues ya asumimos que el trabajo está terminado.
%-XXXXXXXXXXXXXXXXXXXXXXXXXXXXXXXXXXXXXXXXXXXXXXXXXXXXXXXXXXXX

